\subsection{整数指数の場合}

0以上の整数$m,n$に対して
\[
 J_{m,n}=\int_0^1x^m(1-x)^n\dd x
\]
を考える.\,$m,n$が小さければ展開で求められる.しかし指数が大きくなると,展開した項数も多くなる.ここでは$x^m$と$(1-x)^n$をそのまま残し,部分積分で$(1-x)$の指数を1つずつ下げて,\,$n=0$の場合に帰着させることを目標とする.

$n\geqq1$として,\,$x^m$を積分し,\,$(1-x)^n$を微分すると
\begin{align*}
 J_{m,n}
 &=\left[\frac{x^{m+1}}{m+1}(1-x)^n\right]_0^1
   +\frac{n}{m+1}\int_0^1x^{m+1}(1-x)^{n-1}\dd x\\
 &=\frac{n}{m+1}J_{m+1,n-1}
\end{align*}
となる.\,$x=0$では$x^{m+1}$が,\,$x=1$では$n\geqq1$により$(1-x)^n$が$0$になるので,端点の項は消える.

添字が2つあるが,指数の和$m+n$は変わらず,\,$x$の指数が1つ増えて$(1-x)$の指数が1つ減っている.よって$n$が$0$になるまで繰り返せば
\begin{align*}
 J_{m,n}
 &=\frac{n}{m+1}J_{m+1,n-1}\\
 &=\frac{n}{m+1}\cdot\frac{n-1}{m+2}J_{m+2,n-2}\\
 &\qquad\vdots\\
 &=\frac{n}{m+1}\cdot\frac{n-1}{m+2}\cdots\frac{1}{m+n}J_{m+n,0}\\
 &=\frac{n!}{(m+1)(m+2)\cdots(m+n)}J_{m+n,0}
\end{align*}
となる.最後の積分は
\[
 J_{m+n,0}=\int_0^1x^{m+n}\dd x=\frac1{m+n+1}
\]
である.分母は$(m+1)(m+2)\cdots(m+n+1)$となり,分子と分母に$m!$を掛ければ分母は$(m+n+1)!$になるので
\begin{equation}\label{式:整数ベータ積分}
 \boxed{
 \int_0^1x^m(1-x)^n\dd x=\frac{m!\,n!}{(m+n+1)!}}
\end{equation}
を得る.\,$n=0$の場合も,右辺は$\frac{m!}{(m+1)!}=\frac1{m+1}$となり,同じ式で表される.

\begin{bookpoint}
\textbf{POINT}\par
$x^m$を積分し$(1-x)^n$を微分する部分積分で$J_{m,n}=\frac{n}{m+1}J_{m+1,n-1}$とし,\,$n=0$まで下げて$\frac{m!\,n!}{(m+n+1)!}$を得る.
\end{bookpoint}

\begin{example}{ベータ積分}
次の積分を求めよ.
\[
 \int_0^1x^3(1-x)^2\dd x.
\]
\end{example}

\paragraphbf{方針}
式\eqref{式:整数ベータ積分}で$m=3$,\,$n=2$とする.展開しても検算できる.

\paragraphbf{解答}
$\frac1{60}$

\paragraphbf{解法}
式\eqref{式:整数ベータ積分}で$m=3$,\,$n=2$とすると
\[
 \int_0^1x^3(1-x)^2\dd x
 =\frac{3!\,2!}{6!}
 =\frac{12}{720}
 =\frac1{60}
\]
となる.展開で検算すると
\begin{align*}
 \int_0^1(x^3-2x^4+x^5)\dd x
 &=\frac14-\frac25+\frac16\\
 &=\frac1{60}
\end{align*}
となり,一致する.公式があれば展開は要らないが,指数がどこに入るかを確かめるため,小さい場合を計算してみることにも意味がある.
